% !TEX root = ../main.tex
\begin{abstract}
学位论文写作需要把研究问题、材料、方法和结论组织成可追溯的论证。本指南围绕这一过程，介绍章节规划、段落组织、公式与图表表达、引用管理以及提交前自查。指南通过原创对照句说明如何将含糊表述改为边界明确的文字，以三个教学数值展示变量定义和均值计算，并给出逐步修订的流程。对于院系要求和导师意见等尚未确认的信息，指南保留明确的待办标记。示例旨在帮助读者理解可编辑论文工程的组织方式，以及内容审阅、结构检查和真实编译各自能够提供的证据。它不以编译成功代替学术判断，也不把通用建议视为具体学校的提交规范。
\thusetup{keywords={学位论文, 写作方法, 证据管理, 质量检查}}
\end{abstract}

\begin{abstract*}
Thesis writing requires a traceable argument connecting research questions, source material, methods, and conclusions. This guide introduces chapter planning, paragraph development, equations and tables, citation management, and checks before submission. Original sentence pairs illustrate how vague statements can be replaced with bounded descriptions. Three illustrative values demonstrate variable definitions and mean calculation, while a workflow explains iterative revision. Unconfirmed departmental requirements and supervisory preferences remain explicit action items. The example explains the organization of an editable thesis project and the distinct evidence provided by content review, structural checks, and actual compilation. A successful build does not replace academic judgment, and general writing advice does not establish institution-specific submission requirements.
\thusetup{keywords*={thesis writing, writing methods, evidence management, quality checks}}
\end{abstract*}
